problem:  
Let \(G=\mathrm{SO}(p+1)\times \mathrm{SO}(q+1)\) act isometrically on \(S^{p+q+1}\) with cohomogeneity one. The orbit space is \([0,\tfrac{\pi}{2}]\), and each orbit at parameter \(t\) is \(G/H\simeq S^{p}\!\!\times S^{q}\) with metric  
\[
g_t=\cos^2 t\, g_{S^p}+ \sin^2 t\, g_{S^q}.
\]
The well‑known \(G\)‑invariant minimal hypersurface is the Clifford torus \(S^p(\cos t_0)\times S^q(\sin t_0)\) at \(t_0=\arctan\!\sqrt{p/q}\).  

Now perturb the metric within the same cohomogeneity‑one class by  
\[
g_{\varepsilon}=dt^2+\cos^2 t(1+\varepsilon \phi_1(t))^2 g_{S^p}
+\sin^2 t(1+\varepsilon \phi_2(t))^2 g_{S^q},
\]
with smooth even functions \(\phi_i(t)\) preserving the symmetry and smoothness at \(t=0,\tfrac{\pi}{2}\).  

Let \(\mathcal{L}_\varepsilon\) be the \(G\)-invariant Jacobi (stability) operator of the minimal profile curve corresponding to \(S^p(\cos t_0)\times S^q(\sin t_0)\) computed for \(g_\varepsilon\).  

---

**Problem (Equivariant rigidity/bifurcation of the Clifford hypersurface):**  
Does there exist any smooth symmetric perturbation functions \(\phi_1,\phi_2\) and small \(\varepsilon\neq0\) such that the first \(G\)-invariant eigenvalue of \(\mathcal{L}_\varepsilon\) crosses zero, leading (via the Hsiang–Lawson ODE reduction) to a new \(G\)-invariant minimal hypersurface \(S^p_{\varepsilon}\!\times S^q_{\varepsilon}\subset (S^{p+q+1},g_\varepsilon)\) distinct from the deformed Clifford torus?  
Equivalently: are the Clifford hypersurfaces spectrally rigid among all cohomogeneity‑one invariant metric deformations preserving the \(\mathrm{SO}(p{+}1)\times\mathrm{SO}(q{+}1)\) symmetry, or can internal metric warping trigger a genuine equivariant bifurcation of minimal profiles?

---

Why is it a "good" problem:  
This version targets a specific, highly symmetric and elementary situation—the \(\mathrm{SO}(p{+}1)\times\mathrm{SO}(q{+}1)\) action—where all geometric and analytic structures are explicit, yet the outcome is unknown.  It asks whether the classical Clifford minimal tori are rigid not merely in the round sphere but under all metric deformations that preserve their symmetry type.  The question is sharp (spectral crossing of the invariant Jacobi operator), nontrivial (involves singular‑orbit boundary conditions and coupled weight functions in the ODE), and cannot be answered by a general theorem because standard bifurcation theory must adapt to the boundary‑singularity and symmetry restrictions.  

A positive answer would reveal the first example of *intrinsic symmetry‑preserving bifurcation* for minimal submanifolds of a sphere, while a negative answer would establish a new form of **equivariant rigidity** linking the Clifford hypersurface’s spectral data to the spherical geometry.  The problem balances simplicity and depth: it is framed in concrete geometric terms but its resolution demands a synthesis of representation theory, Hsiang–Lawson reduction, and delicate spectral analysis—an archetypal “narrow entrance, deep depth’’ research problem.